\subsection{A uniform nodal norm-thirteen self-isogeny}
\label{subsec:uniform-thirteen-theta-candidate}

\status{Candidate; complete proof and exact certificates under independent audit.}
Work geometrically in characteristic two. Use the already proved normal form
\[
 C_z:\ y^2+y=ax^5+bx^3,\qquad
 a=z^2,\quad b=z^2+z+1,\quad \delta=z+1,
\]
with base point $W$ at infinity. Its normalized geometric modulus satisfies
$\mu^5=(z^2+z+1)^5/z^6$. The order-four automorphism
$g(x,y)=(x+1,y+zx^2+\delta x+\beta)$, where
$\beta^2+\beta=\delta$, is the one constructed in
Lemma~\ref{lem:uniform-f0-order-four-presentation-candidate}.

\begin{lemma}[A cubic algebra for the map $3+2g$]
\label{lem:uniform-thirteen-generic-fibres-candidate}
\status{Candidate; complete exact proof.}
Suppose $z\ne0,1$. Put $\psi=3+2g$, an endomorphism of Rosati norm
$13$, and, at $P=(p,y_0)$, put
\[
\begin{aligned}
 d_1&=ap^4+bp^2,&d_2&=d_1^2+bp,\\
 G&=zp^2+\delta p+\beta,&F&=G+d_1,\\
 Z&=\delta+F+d_2,&
 A_1&=\delta d_1,\qquad A_2=Zd_1+\delta d_2,\\
 A_3&=Z(F+\delta)+\delta^2.
\end{aligned}
\tag{T13.1}
Then $\ker\psi\cap\Theta=\{W\}$. A fibre with at least three
distinct points and $Z\ne0$ satisfies one of the following conditions
for some finite $L$:
\[
 P_1(L)=P_2(L)=0,\qquad\text{or}\qquad A(L)=T(L)=0,
\tag{T13.2}
\]
where the polynomials over $\mathbf F_2(z)$ are defined by
\[
\begin{aligned}
 e&=\delta^{-2},&j&=aL+\delta^2,&k&=\delta^4,\\
 A&=a^2L^7+(a^2+1)L^3+(a+1)L+a,\\
 H&=a^2L^6+(a^2+1)L^2+(a+1),&J&=a^2L^4+a^2+1,\\
 B_2&=Je^2+a^2e^4j^2,&B_1&=1+He+a^2e^4k,\\
 P_2&=B_2^2j^2+B_1^2+LA B_2,&
 P_1&=B_2^2k+LA B_1+A^2e,\\
 T&=B_1^3+j^2B_1B_2^2+kB_2^3.
\end{aligned}
\tag{T13.3}
\]
Fibres involving $W$ or $Z=0$ have at most two points.
\end{lemma}
\begin{proof}
Write $X=x-p$. An arbitrary section of $\mathcal O(7W)$ vanishing
to order at least three at $P$ has the form
\[
 (lX+c)(y+y_0)+cd_1X+(ld_1+cd_2)X^2+qX^3.
\]
Vanishing at $gP$ determines $q=(l+c)F+cd_2$. The remaining
derivative condition there is exactly $\delta l+Zc=0$.
Since $\delta\ne0$, this solution space is one-dimensional at every
finite $P$. Its generator is
\[
 f_P=(ZX+\delta)(y+y_0)+A_1X+A_2X^2+A_3X^3,
\]
with divisor $3P+2gP+E_P-7W$ for a unique effective nonspecial
degree-two divisor $E_P$. Indeed its section space has dimension
one; a special degree-two divisor on this genus-two curve would
have two sections. If $\psi(P-W)=0$, the same section space would
be $H^0(2W)$, a contradiction. This also proves the assertion about
the kernel. Integer coefficients in $\psi$ are endomorphism
coefficients: $\psi^\dagger\psi=(3-2g)(3+2g)=[13]$.

Direct polynomial identities give
\[
 F^2=Z+ap,\qquad
 Z=a^2p^8+(a^2+1)p^4+(a+1)p^2+ap+\beta+\delta,
\tag{T13.4}
\]
and the norm of $f_P$ is
\[
 X^3(X+1)^2
 \bigl(aZ^2X^2+(A_3^2+apZ^2)X+\delta Z^2+a\delta^2\bigr).
\tag{T13.5}
\]
At $Z=0$ one has $A_3=\delta^2\ne0$ and $E_P=W+R$, where
$x(R)=p+a/\delta^2$. Thus a fixed $R$ determines $p$ uniquely.
Two points with the same $p$ are hyperelliptic conjugates. Their
images under $\psi$ could coincide only at a geometric two-torsion
point; in two-rank zero this is zero, already excluded for finite
$P$. Hence the $Z=0$ fibres are harmless. They cannot meet a
$Z\ne0$ fibre: the unique nonspecial residual divisor contains
$W$ in the first case and is finite in the second. Infinity is the
only point over the zero image.

For $Z\ne0$, write the residual divisor in global monic form
$x^2+Lx+M$. Equation (T13.5) gives
\[
 aL=Z+\delta^2+\delta^4/Z^2,\qquad
 p^2+Lp=M+\delta/a+\delta^2/Z^2.
\tag{T13.6}
\]
Consequently $Z$ is a root of $Z^3+jZ^2+k$, a separable cubic:
its derivative is $Z^2$ and $k\ne0$. Set
$U=M+\delta/a+ej$. In this cubic algebra,
$\delta^2/Z^2=e(Z+j)$, and hence $p^2+Lp=U+eZ$.
Reducing (T13.4) by these two relations gives
\[
 A(L)p=B_2Z^2+B_1Z+B_0,
 \quad B_0=HU+JU^2+a^2U^4+\beta+\delta+a^2e^4jk.
\tag{T13.7}
\]
If $A\ne0$, each $Z$ determines $p$ uniquely. A fibre with three
points therefore uses the three distinct cubic roots. Substitution
of (T13.7) in $p^2+Lp=U+eZ$ gives a polynomial of degree at most
two in $Z$; its $Z^2,Z$ coefficients are exactly $P_2,P_1$.
Both must vanish, proving the first condition in (T13.2).

If $A=0$ and the three values $B_2Z^2+B_1Z$ are distinct, a fixed
$U$ selects at most one $Z$, and the quadratic equation for $p$
gives at most two points. If $B_2\ne0$, equality at two cubic roots
forces their sum to be $B_1/B_2$; the third root is $j+B_1/B_2$.
Substitution in the cubic, multiplied by $B_2^3$, is $T=0$.
For $B_2=0$ the three values are distinct unless $B_1=0$, and
$T=B_1^3$ covers exactly that remaining possibility. These are
necessary conditions only; no converse is asserted.
\end{proof}

\begin{lemma}[All moduli except one have a good $3+2g$ presentation]
\label{lem:uniform-thirteen-generic-parameters-candidate}
\status{Candidate; exact polynomial certificate.}
For every normalized modulus with $\mu^5\ne1$, there is a presentation
above for which neither condition (T13.2) occurs.
\end{lemma}
\begin{proof}
Clear the powers of $z+1$ in $P_1,P_2,T$ as recorded by their exact
coefficient arrays. The two resultants
\[
 R_1(z)=\operatorname{Res}_L(\widetilde P_1,\widetilde P_2),
 \qquad R_2(z)=\operatorname{Res}_L(A,\widetilde T)
\tag{T13.8}
\]
are nonzero of degrees $224,198$. Take the product $B_z$ of the
distinct displayed factors of these two polynomials, also including
$z(z+1)$. It has degree $137$. The complete coefficients and
factor-product identities are retained; factor irreducibility is
unnecessary for this zero-set implication.

Put $r=\xi^2+\xi$, $z=(r^2+1)/r^3$, and remove powers of
$H_4(\xi)=\xi^4+\xi$ from the homogeneous pullback of $B_z$.
The resulting polynomial $B$ has degree $816$. Write
$B=B_0+B_1+B_2$ by exponent classes modulo three. The exact identity is
\[
\gcd\bigl(B(V),B(V+1),V^{816}B(1/V),B_0+B_2,B_1+B_2\bigr)
       =D_{12}(V)^5+H_4(V)^{12},
\tag{T13.9}
\]
where $D_{12}=V^{12}+V^9+V^6+V^3+1$.
The right side is also $\sum_{i=0}^{20}V^{3i}$.
All identities here are in $\mathbf F_2[V]$ and have reconstructive
coefficient and Bezout certificates.

The modulus equation is
$\mu^5=D_{12}(\xi)^5/H_4(\xi)^{12}$, as in the eleven-case
parameter proof. Its monic degree-sixty numerator has a root
outside $\mathbf F_4$, since $D_{12}H_4=V^{16}+V$ and $D_{12}$
is nonzero at every zero of $H_4$. If all five presentations
$\xi,\xi+1,\xi^{-1},\omega\xi,\omega^2\xi$ were bad, where
$\omega^2+\omega+1=0$, (T13.9) would force $\mu^5=1$.
Thus one is good. The usual fifth-root normalization gives the
prescribed $\mu$, not merely its fifth power. This proves the lemma,
including $\mu=0$.
\end{proof}

\begin{lemma}[A different norm-thirteen map at the exceptional modulus]
\label{lem:uniform-thirteen-exceptional-map-candidate}
\status{Candidate; complete finite-field certificate and geometric proof.}
The geometric isomorphism class $\mu^5=1$ has an etale polarized
thirteen-self-isogeny whose restriction to its principal theta has
no fibre with three distinct points.
\end{lemma}
\begin{proof}
The following exact finite field fixes all coefficient inputs:
\[
 k_0=\mathbf F_2[c]/
 (c^{12}+c^7+c^6+c^5+c^3+c+1).
\tag{T13.10}
\]
Encode $\sum e_ic^i$ by the integer $\sum e_i2^i$.
Take $z=1971$ in this encoding; it satisfies $z^3+z^2+1=0$,
so the corresponding $C_z$ has $\mu^5=1$.
For $i=1,2,3$, let
\[
 g_i(x,y)=(x+\tau_i,y+u_ix^2+v_ix+\beta_i),\qquad
 (\tau_i,u_i,v_i,\beta_i)=
 \begin{cases}
 (1,1971,1970,3526),\\
 (970,2395,1,2854),\\
 (1145,3817,1,3010).
 \end{cases}
\tag{T13.11}
\]
Direct coefficient checks establish that these preserve the curve,
fix $W$, square to the hyperelliptic involution, and pairwise
anticommute on its Jacobian. Explicitly the checks are
\[
\begin{gathered}
 u_i^2=a\tau_i,\quad v_i^2+u_i=b\tau_i,\quad
 v_i=a\tau_i^4+b\tau_i^2,\quad
 \beta_i^2+\beta_i=a\tau_i^5+b\tau_i^3,\\
 u_i\tau_i^2+v_i\tau_i=1,\quad
 u_i\tau_j^2+v_i\tau_j+u_j\tau_i^2+v_j\tau_i=1\quad(i\ne j).
\end{gathered}
\tag{T13.12}
\]
Thus $\phi=1+2g_1+2g_2+2g_3$ satisfies
$\phi^\dagger\phi=[1+4+4+4]=[13]$.
It is etale of degree $169$. The three translation classes are
$\mathbf F_2$-independent; this construction is not the earlier
single-quaternion expression $1+2g+2h+2gh$.

Here is a fully specified interpolation procedure, and the exact
certificates needed to control every fibre. Write $X=x-p$ and
$Y=y+y_0$. A section of $\mathcal O(9W)$ vanishing at $P$ is
\[
 f=(l_0+l_1X+l_2X^2)Y+c_1X+c_2X^2+c_3X^3+c_4X^4.
\tag{T13.13}
\]
For $s_i=\tau_i$, put $Y_i=u_ip^2+v_ip+\beta_i$ and
$D_i=a(p+s_i)^4+b(p+s_i)^2$. For each of the three coordinate
vectors $l=(l_0,l_1,l_2)$, set $l(s)=l_0+l_1s+l_2s^2$ and
\[
\begin{aligned}
 F_i&=l_1Y_i+l(s_i)D_i,\\
 c_3&=(F_1+F_2)/(s_1^2+s_2^2),\quad c_1=F_1+c_3s_1^2,\\
 V_i&=l(s_i)Y_i+c_1s_i+c_3s_i^3\quad(i=1,2),\\
 c_4&=(V_2s_1^2+V_1s_2^2)/(s_1^2s_2^4+s_1^4s_2^2),\quad
 c_2=(V_1+c_4s_1^4)/s_1^2.
\end{aligned}
\tag{T13.14}
\]
All denominators are nonzero constants. The remaining derivative
and value equations at $g_3P$ form a two-by-three matrix over
$k_0[p]$. Its vector of maximal minors gives the three $l_i$;
(T13.14) then gives the $c_i$. The gcd of the three minors is one,
verified with a polynomial Bezout relation. Consequently the matrix
has rank two at every finite $p$, and the seven imposed conditions
have a one-dimensional solution space. This proves that the residual
degree-two divisor $E_P$ is always nonspecial and unique, and that
$\ker\phi\cap\Theta=\{W\}$, by the same Riemann--Roch argument
as above.

The exact norm factorization of this section is
\[
 f\iota(f)=X\prod_{i=1}^3(X+s_i)^2
                 \bigl(A(p)X^2+B(p)X+C(p)\bigr),
 \qquad A=a l_2^2.
\tag{T13.15}
\]
The degrees of $A,B,C$ are $16,24,16$. Every coefficient is
reconstructed from (T13.10)--(T13.14); no sampled curve equation
or inferred numerical pattern is used. On $l_2\ne0$, the residual
monic quadratic in global $x$ has coefficients
\[
 \mathcal L=B/A,\qquad
 \mathcal M=p^2+pB/A+C/A.
\tag{T13.16}
\]
Their reduced numerator/denominator degrees are respectively
$(24,16)$ and $(25,16)$, and their denominator zero sets are
exactly the zeros of $l_2$.

For either rational function $f=N_f/D_f$ in (T13.16), define
\[
 E_f(p,q)=
 \frac{N_f(q)D_f(p)+N_f(p)D_f(q)}{q+p}\in k_0[p,q].
\tag{T13.17}
\]
The division is exact. The retained coefficient certificate provides
explicit polynomials $U,V\in k_0[p,q]$ and $S_0,S_1\in k_0[p]$
such that
\[
 UE_{\mathcal L}+VE_{\mathcal M}=S_0(p)+S_1(p)q,
 \qquad \gcd(S_0,S_1)=\overline l_2^{\,44},
\tag{T13.18}
\]
where $\overline l_2$ is the monic normalization of $l_2$.
The degrees of $S_0,S_1$ are $877,876$, and the $q$-degrees of
$U,V$ are $22,21$. The complete arrays are retained, and a direct
polynomial multiplication verifies (T13.18); this proof does not
invoke specialization of a generic polynomial gcd. At a fixed
$p$ with $l_2(p)\ne0$, the right side is a nonzero polynomial
of degree at most one in $q$. Two other distinct nonpole points
with the same residual divisor would give two distinct roots,
which is impossible. Thus every nonpole fibre has at most two
distinct $x$-coordinates.

For completeness the poles are checked separately. The polynomial
$l_2$ is squarefree of degree eight and is relatively prime to
$c_4$; modulo $l_2$, one has $B=c_4^2\ne0$. Hence at these
eight points the section has pole order eight and $E_P=W+R$,
with
\[
                    x(R)=p+C/B\pmod {l_2}.
\tag{T13.19}
\]
The eight vectors $1,x(R),\ldots,x(R)^7$ are a basis of
$k_0[p]/(l_2)$. The minimal polynomial of $x(R)$, in the encoding
(T13.10) and in ascending coefficient order, is
\[
 (2956,4012,3611,1,3655,2676,2677,0,1).
\tag{T13.20}
\]
It is separable. Therefore the eight pole points have pairwise
distinct residual-point $x$-coordinates. Pole residual divisors
cannot coincide with nonpole ones, since the former contain $W$
and the latter are finite. Finally, equal source $x$-coordinates
would give hyperelliptic conjugates with equal image, hence a
geometric two-torsion image, necessarily zero in two-rank zero.
The already proved kernel assertion excludes this except at $W$.
These arguments cover infinity, all poles, repeated residual roots,
and all finite points, and establish the lemma.
\end{proof}

\begin{theorem}[A nodal thirteen-self-isogeny on every two-rank-zero Jacobian]
\label{thm:uniform-thirteen-nodal-self-isogeny-candidate}
\status{Candidate; assembly of the preceding exact proofs.}
Every indecomposable two-rank-zero principally polarized abelian surface
in characteristic two has an etale polarized self-isogeny of multiplier
$13$ whose principal theta image has exactly $168$ ordinary nodes
and no other singularities.
\end{theorem}
\begin{proof}
For $\mu^5\ne1$ use the first two lemmas. For $\mu^5=1$ use
the explicitly different map in the third lemma. The normal-form
classification identifies every such exceptional surface with its
specified finite-field model after geometric base extension.
In both cases every theta fibre has at most two points. Apply
Lemma~\ref{lem:root-nonordinary-odd-theta-nodality}: the image is
nodal, with $|\ker\phi|-1=13^2-1=168$ nodes.
\end{proof}

\begin{corollary}[All odd degrees divisible by thirteen in two-rank zero]
\label{cor:f0-odd-multiples-thirteen-candidate}
\status{Candidate; complete group-theoretic transfer from the preceding theorem.}
On every fixed two-rank-zero principally polarized surface $(A,\Theta)$
over a perfect field of characteristic two, put
$L=\mathcal O_A(3\Theta)$ and use the complete systems
$|L^d|=|3d\Theta|$. Every positive odd $d$ divisible by at least one
of $3,5,7,11,13$ has relative geometric and arithmetic image
$\operatorname{GL}_{9d^2-1}(\mathbf F_2)$. Every even degree has
the previously proved full one-covector stabilizer. The first
remaining odd fixed-surface degree in two-rank zero is $17$.
\end{corollary}
\begin{proof}
The good-theta norm step, Theorem~\ref{thm:good-theta-odd-norm-step-candidate},
gives degree thirteen from degree one. For $d=13n$, choose a
nonproduct incoming source of odd multiplier $n$ and use
Theorem~\ref{thm:cubic-component-arbitrary-odd-norm-transfer-candidate}
with $e=13$ when $n>1$. The case $n=1$ is the preceding step.
Products and the other listed degrees were already proved. Full
geometric general linear image forces the same arithmetic image;
the geometric isogeny need not be defined over the original field.
\end{proof}

\begin{remark}[Verification scope and retained failed choices]
\label{rem:uniform-thirteen-verification-scope-candidate}
\status{Candidate; exact computation and audit record.}
The universal interpolation and parameter programs are
\path{code/monodromy/post_CP28_thirteen_interpolation.py} and
\path{code/monodromy/post_CP28_thirteen_parameter_probe.py}.
The exceptional interpolation, coordinate, pole and polynomial-combination
programs have stems \path{post_CP28_thirteen_three_anticommuting_probe},
\path{post_CP28_thirteen_three_unit_coordinates},
\path{post_CP28_thirteen_pole_fibres}, and
\path{post_CP28_thirteen_linear_subresultants} in the same exported
directory. Their complete coefficient outputs are retained under
\path{agents/monodromy/outputs/}.
The compact reconstructive checker is
\path{code/monodromy/verify_nonordinary_thirteen_certificates.py}.
The compact checker passed all 34 exact checks at 00:15:05 UTC on
20 November 2026 in 0.723901527 seconds. It independently rebuilds
both generic resultants by binary-polynomial Bareiss determinants,
checks the field modulus by the degree-twelve irreducibility test,
and verifies the raw-minor Bezout identity, the entire stored
polynomial combination, and its exact coefficient gcd.
The exact $(\mathcal L,\mathcal M)$ combination alone proves (T13.18);
a second, redundant $(\mathcal L,\nu)$ combination was computed
and is retained but is not needed here.

The original choice $3+2g$ really can fail at $\mu^5=1$:
an independently reconstructed four-point packet is retained.
The alternative $1+2g+2h+2gh$ in a single quaternion subgroup also
has interpolation rank drops in the retained test. Neither is used
for the exceptional surface in this theorem. Two slower exploratory
remainder-sequence runs timed out while factoring unnecessary
exception polynomials; their code and outputs remain preserved.
The final argument instead checks the explicit polynomial identity
(T13.18) and its coefficient gcd directly. No sampling, unproved
resultant-specialization inference, or higher-degree Cartier model
is used.
\end{remark}
